\begin{proof}
Let  

\[
\mathbf a=(a_{1},\dots ,a_{n}), \qquad 
\mathbf b=(b_{1},\dots ,b_{n})\in\mathbb{R}^{n},
\]

and denote the standard Euclidean inner product and norm by  

\[
\langle\mathbf a,\mathbf b\rangle =\sum_{i=1}^{n}a_{i}b_{i},
\qquad 
\|\mathbf a\| =\sqrt{\langle\mathbf a,\mathbf a\rangle}
           =\sqrt{\sum_{i=1}^{n}a_{i}^{2}} .
\]

For any real scalar \(t\) consider the non‑negative quadratic form  

\[
f(t)=\|\mathbf a+t\mathbf b\|^{2}
     =\sum_{i=1}^{n}(a_{i}+t b_{i})^{2}\ge 0\qquad (t\in\mathbb R).
\]

Expanding the square and collecting terms gives  

\[
\begin{aligned}
f(t)
&=\sum_{i=1}^{n}\bigl(a_{i}^{2}+2t\,a_{i}b_{i}+t^{2}b_{i}^{2}\bigr)  \\
&=\underbrace{\sum_{i=1}^{n}a_{i}^{2}}_{A}
   +2t\underbrace{\sum_{i=1}^{n}a_{i}b_{i}}_{B}
   +t^{2}\underbrace{\sum_{i=1}^{n}b_{i}^{2}}_{C}  \\
&=A+2Bt+Ct^{2}.
\end{aligned}
\]

Thus \(q(t):=Ct^{2}+2Bt+A\) satisfies \(q(t)=f(t)\ge0\) for every \(t\in\mathbb R\).

\medskip
\noindent\textbf{Case \(C>0\).}  
Since the leading coefficient \(C\) is positive, a quadratic is non‑negative on \(\mathbb R\) iff its discriminant is non‑positive:
\[
(2B)^{2}-4AC\le 0.
\]
Dividing by \(4>0\) yields
\[
B^{2}\le AC.
\]

\noindent\textbf{Case \(C=0\).}  
Then \(\mathbf b=\mathbf0\), whence \(B=0\) and the desired inequality reduces to \(0\le0\), which holds.

Combining the two cases we obtain
\[
\biggl(\sum_{i=1}^{n}a_{i}b_{i}\biggr)^{2}
\le
\biggl(\sum_{i=1}^{n}a_{i}^{2}\biggr)
\biggl(\sum_{i=1}^{n}b_{i}^{2}\biggr),
\]
i.e. the Cauchy–Schwarz inequality.

\medskip
\textbf{Equality case.}  
Equality in the discriminant condition occurs exactly when  

\[
(2B)^{2}-4AC=0\;\Longleftrightarrow\;B^{2}=AC .
\]

If \(C>0\), this means that \(q(t)\) has a double root  
\(t_{0}=-\dfrac{B}{C}\), and therefore  

\[
\mathbf a+t_{0}\mathbf b=\mathbf0 \quad\Longrightarrow\quad
\mathbf a=-t_{0}\mathbf b .
\]

Thus the vectors are linearly dependent: there exists \(\lambda\in\mathbb R\) such that
\(\mathbf a=\lambda\mathbf b\), i.e. \(a_{i}=\lambda b_{i}\) for all \(i\).
Conversely, if \(\mathbf a=\lambda\mathbf b\) for some \(\lambda\), then both sides of the
inequality are equal. Hence equality holds iff the two vectors are proportional.

\end{proof}